\section{Discussion}
\label{sec:discussion}

The central result is not a single optimal parameter combination but a measurement distinction. Source popularity and misinformation dissemination can move in opposite directions. Full copying produced the largest increase in target-source selection, yet it reduced realised fake reposting because it transferred credibility and thereby lowered the target's probability of producing false content. Non-credibility copying produced a smaller popularity gain but a larger misinformation increase because it preserved the source's editorial propensity. Studies that label every selection of a ``fake-news source'' as a fake repost would miss this difference and overstate the harm of full copying while potentially understating camouflage.

This finding has a direct modelling implication. Perceived credibility and objective false-publication probability should ideally be distinct state variables. Their coupling in the current implementation made the diagnostic possible, but it also caused one scenario parameter to change both user perception and source behaviour. Separating them would support interventions such as improving reputation without changing editorial conduct, or changing editorial conduct without immediate reputation repair. Until that refactoring is made, non-credibility copying is the cleanest available dissemination treatment and credibility/all-copy scenarios should be labelled mechanistic diagnostics.

The baseline supports the internal logic of truth-sensitive WOM. Discovery alone did not alter the retained outcomes, whereas penalising false and rewarding true recommendations reduced fake reposting substantially. This is consistent with research showing that attention to accuracy and source quality can improve information choice \citep{pennycook2019,pennycook2020,epstein2021}. The comparison remains model specific: it does not estimate the effect of a platform fact-check or a real interpersonal correction. It shows that signed truth feedback, when integrated into an endorsement-based decision process, can shift aggregate repost composition.

The confirmation results expose why policy labels must be interpreted relationally. The largest camouflage-control difference occurred under false penalisation and no true reward. At first glance this seems to conflict with the protective baseline. It does not. The baseline compares policies, while confirmation compares scenarios within a policy. Under $-1/+1$, both treatment and control benefit from positive evidence attached to true outcomes, and their late-run difference is small. Under $-1/0$, that stabilising pathway is absent; copied engagement attributes create a much larger relative separation even though false outcomes are penalised in both branches. Future work should decompose the endorsement stock over time to determine precisely how positive true-news WOM, negative false-news WOM, and non-WOM source attributes contribute to the score trajectories.

Memory and timing behave as coupled path-dependence mechanisms. A very short memory removes scenario-generated evidence before it can accumulate, while an early intervention under long memory shapes a large fraction of the run. Delayed interventions have less time to alter the evidence distribution and must compete with pre-existing endorsements. These patterns resemble cumulative-advantage processes in which early social signals can produce persistent inequality \citep{salganik2006,muchnik2013}. However, the model does not imply that longer human memory necessarily increases vulnerability. Infinite model memory preserves every form of evidence, including negative evidence and control history. Its effect depends on what evidence the scenario and WOM policy generate.

The two network analyses illustrate a broader lesson about robustness. An effect is meaningful only relative to its control. WOM-on/off contrasts show that truth-sensitive social exchange can suppress the transformed source across supported contact and reach cells. Camouflage/control contrasts show that copied attributes can still add a small burden after WOM is enabled. Reporting either result without its estimand would create an apparent contradiction. For platform design, the distinction corresponds to asking whether social feedback is beneficial overall versus whether a manipulation remains effective after social feedback exists.

The common-seed analysis improves confidence in the ordering of conditions. Near-perfect rank reproduction indicates that the major parameter patterns are not artefacts of a particular independent Monte Carlo sample. Pairing usually narrows uncertainty, but not automatically. Treatment branches change later random calls, and nearly zero effects can have unstable standard-error ratios. The experiment harness therefore records seeds and reports empirical precision gains rather than assuming them.

\subsection{Implications for research and practice}

For misinformation research, the study supports a three-part outcome framework: audience share, editorial false-publication propensity, and realised false reposts. These quantities answer different questions and should be reported together. For journalism, a source that adopts familiar professional conventions may gain legitimacy without changing its truthfulness. For communication and marketing research, engagement attributes can transfer attention, but the effect depends on when they are introduced and on the audience's retained evidence. For platform engineering, recommendation exposure, outcome evaluation, and source reach should be logged as separate mechanisms.

The model may also support countermeasure experiments. Media-literacy interventions can be represented by changing user weights on credibility, information quality, sensationalism, and emotion. Platform friction can be represented by reduced reach or delayed transmission. Reputation systems can be represented by truth-sensitive WOM with different magnitudes. Yet such scenarios should be tested only after separating perceived credibility from objective false-publication probability, because otherwise an intervention that changes perceived trust may accidentally rewrite source behaviour.

\subsection{Limitations}

First, the model is not empirically calibrated. The parameter values originate from the project's input workbook and earlier modelling design, not a joint fit to platform exposure and repost data. Monte Carlo intervals describe stochastic variation conditional on those values; they are not population confidence intervals. The repository includes a calibration template and weighted-distance tool, but no empirical targets, sampling uncertainties, or held-out observations.

Second, the network representation is limited. Contact opportunity and a binary source-reach mechanism are supported, but homophily is absent. Echo-chamber dynamics documented in empirical research \citep{delvicario2016} cannot be evaluated by the current CLI harness. Degree is also represented through a compact contact mechanism rather than a validated platform network topology.

Third, sources are types rather than evolving organisations. They do not strategically choose truth states, react to moderation, or adapt attributes in response to audience feedback. The scenario is exogenous and deterministic once its activation period is reached. Users likewise do not create original content, verify claims independently, or leave the platform.

Fourth, the factorial strategy ranking involves 108 effects and no correction for multiple comparisons. We therefore treat rankings as exploratory and emphasise the 30-seed confirmation for policy claims. Even there, the six conditions were selected after inspecting earlier patterns, so a future preregistered replication would provide stronger confirmatory evidence.

Finally, common random numbers do not guarantee paired trajectories. Once a treatment changes source selection, subsequent random-number consumption can diverge. Pairing remains useful empirically, but the observed precision improvement should not be interpreted as perfect counterfactual synchronisation.

\subsection{Next research steps}

The next technical step is to separate perceived credibility, reputation, and objective false-publication probability. The next empirical step is to obtain targets for source exposure, repost rates conditional on exposure, source-level truthfulness, and credibility judgments. Calibration should estimate WOM weight, memory horizon, reach, and the endorsement aggregation scale using training targets, then evaluate predictive ordering on held-out targets. Finally, a homophily mechanism and alternative network topologies should be introduced through an explicit model extension rather than simulated indirectly by workbook manipulation.
